\chapter{LITERATURE SURVEY}

The review covers the study that produced the dataset, work that benchmarks
classifier families on the same features, work that addresses the class
imbalance, and recent attempts that add preprocessing or deep architectures.
Methodological sources for the individual algorithms are cited in Chapter~5.

\section{Prior Work}

\subsection{The originating study}

Köklü and Özkan \cite{koklu2020} built a computer vision system to separate seven
registered Turkish dry bean varieties. Images of 13,611 grains were segmented and
reduced to 16 descriptors, 12 dimensional and 4 shape factors. Four classifiers
were trained with 10-fold cross-validation: multilayer perceptron (MLP), support
vector machine (SVM), k-nearest neighbours (kNN) and decision tree (DT). The
reported correct classification rates were 93.13\% for SVM, 92.52\% for DT,
91.73\% for MLP and 87.92\% for kNN, and a MATLAB interface was built around the
best model. The study establishes that morphological features alone are sufficient
for high accuracy on these varieties, and it released the feature table that makes
replication possible.

Two limitations matter here. The comparison is reported as a single accuracy per
model, which is uninformative when the smallest class has 522 beans against 3,546
for the largest: a model can raise accuracy by favouring Dermason. And the table
contains strongly collinear descriptors, without any examination of whether all 16
are needed. Both are addressed here through macro-averaged metrics and a feature
ablation.

\subsection{Benchmarking SVM kernels on the same dataset}

Mehta et al. \cite{mehta2023} compared linear, polynomial and radial basis function
SVM kernels with other popular classifiers on the same dataset, applying principal
component analysis as preprocessing and using accuracy as the primary metric. The
RBF kernel gave the best result, with an accuracy of 93.34\%, precision of 92.61\%,
recall of 92.35\% and F1-score of 91.40\%, and the conclusion is that a nonlinear
kernel is appropriate because the class boundaries are not hyperplanes.

The study isolates the kernel choice, which is one of the decisions taken here. Its
limitation is that PCA is applied without a nested protocol and only accuracy is
reported, so the effect of dimensionality reduction on the minority class is not
visible, and a single train/test configuration gives no sense of the variance.

\subsection{Class imbalance and hybrid resampling}

Lee, Wang and Huang \cite{lee2024} argued that the dry bean data are imbalanced and
that plain DT, RF and SVM therefore under-perform on the minority variety. They
combined borderline-SMOTE oversampling with k-means clustering and trained DT, RF
and SVM on the clustered, balanced data with 10-fold cross-validation. The plain
models reached test accuracies of 88.83\% (DT), 83.25\% (RF) and 93.75\% (SVM),
while BLSMOTE + k-means + SVM reached 97.54\% with a precision of 0.9736 and an AUC
of 0.9831.

This is the strongest published accuracy on the dataset, but the protocol applies
oversampling before the split, so synthetic minority samples generated from a bean
that also appears in the test partition can inflate the measured accuracy. The
present study therefore evaluates the original distribution first, reports macro
F1, and tests class weighting as an explicit leakage-free alternative. The result
in Chapter~6 is a useful negative finding rather than a contradiction of
\cite{lee2024}.

\subsection{Phenotyping, dimensionality reduction and deep alternatives}

Mendigoria et al. \cite{mendigoria2021} extended the use of these descriptors
beyond the class label, combining image-derived seed traits with supervised
learning to predict architectural phenes of dry bean seed as well as variety. The
motivation is agronomic: identifying seed characteristics early supports quality
control and selection. The study confirms that the 16 variables carry information
about more than the label, which is relevant to the interpretation in Chapter~6,
where importance is discussed in terms of bean geometry.

Gebeş and Şahin \cite{gebes2025} compared several classifiers on dry bean data and
measured how different dimensionality reduction techniques change classification
success. Their framing matches the question raised by the correlation analysis in
Chapter~3: several descriptors are near-duplicates, so compression might be
expected to help. Their work motivates the ablation in Section~6.4, which tests the
same question with cross-validated macro F1 instead of accuracy.

Phan \cite{phan2026} proposed a lightweight one-dimensional residual network that
treats the 16 descriptors as a short sequence rather than as independent columns.
A 1D ResNet can model interactions between neighbouring descriptors, but it also
illustrates the cost side of the trade-off: the tabular problem that a tuned
RBF-SVM solves in seconds is being solved with a network that needs more data, more
tuning and more compute. This case study stays with classical learners and
Chapter~6 quantifies what that costs in accuracy.

\section{Research Gap and Motivation}

Three gaps follow. Published comparisons report accuracy, or a single aggregate F1,
on an imbalanced seven-class problem where the smallest variety contributes 3.8\%
of the beans; macro-averaged precision, recall and F1 are needed to show whether a
model has learned Bombay or has effectively dropped it. The published protocols
also differ in ways that make the numbers incomparable: some use 10-fold
cross-validation on the full dataset, some a single split, and at least one applies
oversampling before the split. And the cost of each model is rarely reported,
although it is a practical selection criterion for a laptop-scale workflow.

This case study addresses those points on a single, fixed, stratified split; it
selects models by cross-validated macro F1 rather than accuracy; it reports
per-class behaviour, confusion structure and computation time alongside the
aggregate metrics; and it tests the two preprocessing questions raised in the
literature, redundancy and class weighting, as controlled experiments rather than
assumptions. No new algorithm is proposed, and no claim of state-of-the-art
accuracy is made: the contribution is a transparent, reproducible comparison.

\begin{table}[htbp]
\centering
\caption{Comparison of the reviewed studies}
\label{tab:literature}
\footnotesize
\renewcommand{\arraystretch}{1.12}
\begin{tabularx}{\textwidth}{@{}>{\raggedright\arraybackslash}p{2.05cm}>{\raggedright\arraybackslash}p{2.5cm}>{\raggedright\arraybackslash}p{2.1cm}>{\raggedright\arraybackslash}p{2.6cm}>{\raggedright\arraybackslash}X@{}}
\toprule
\textbf{Study} & \textbf{Objective} & \textbf{Data} & \textbf{Methods} & \textbf{Findings and limitations} \\
\midrule
Köklü and Özkan \cite{koklu2020} &
Separate seven registered varieties with a vision system &
13,611 beans, 16 features (dataset source) &
MLP, SVM, kNN, DT; 10-fold CV &
SVM best at 93.13\%; single accuracy figure, no per-class analysis, no feature ablation \\
Mehta et al. \cite{mehta2023} &
Benchmark SVM kernels and classifiers &
Same dataset, PCA preprocessing &
Linear, polynomial and RBF SVM; other classifiers &
RBF SVM 93.34\% accuracy, F1 91.40\%; accuracy-only metric, single split \\
Lee et al. \cite{lee2024} &
Fix imbalance before classification &
Same dataset, 80/20 split &
BLSMOTE + k-means with DT, RF, SVM &
97.54\% test accuracy with BLSMOTE + k-means + SVM; oversampling applied before splitting, inflating results \\
Mendigoria et al. \cite{mendigoria2021} &
Predict seed architectural phenes and classify variety &
Image-derived seed traits &
Supervised learners on morphometric traits &
Confirms descriptors carry trait information beyond the label; abstract-level detail only \\
Gebeş and Şahin \cite{gebes2025} &
Measure the effect of dimensionality reduction &
Dry bean feature table &
Multiple classifiers with and without reduction &
Reduction changes classification success; motivates the ablation in Section 6.4 \\
Phan \cite{phan2026} &
Improve accuracy with a deep sequence model &
Dry bean feature table &
Lightweight 1D-ResNet &
Shows descriptor interactions can be learned; higher tuning and compute cost than classical learners \\
\bottomrule
\end{tabularx}
\end{table}
